\scanpage{46}
% archive sheet number: 45
Lorsque la $n$-carte $\Pi$ est un polyèdre (= géométrie d'incidences), alors la donnée de $\unsure{\varphi|\phi}_R$ satisfaisant (2) \uline{équivaut} à la donnée des $\varphi_j$ $(0\leq j\leq n)$, \unsure{soumises|soumis} à la condition d'être compatibles avec les relations d'incidences. Dans le cas général, on peut dire \inserted{encore} que (3) se généralise en des applications

\diagram{\begin{tikzcd}[ampersand replacement=\&,column sep=huge]
(4)\quad \unsure{D|\mathscr{D}}_{\{j_0,j_1,\ldots,j_p\}}(\Pi) \arrow[r,"\varphi_{\{j_0,\ldots,j_p\}}"] \& \operatorname{Drap\ aff}_{\{j_0,j_1,\ldots,j_p\}}(\Pi)
\end{tikzcd}}

pour \unsure{tte} partie $\{j_0<j_1<\cdots<j_p\}$ de $\Delta_n:\{0,1,\ldots,n\}$, ces applications étant compatibles avec les applications verticales pour des inclusions de parties de $\Delta_n$. La donnée de $\varphi_R$ satisfaisant (2) équivaut à la donnée des applications (4) satisfaisant la condition de fonctorialité précédente --- on peut donc les interpréter \inserted{les $\varphi_R$} comme des hom. de géométries de drapeaux basées sur \inserted{l'ens\unsure{ble} ordonné} $\Delta_n$, de $D_{\unsure{*|\bullet}}(\Pi)$ dans $\operatorname{Drap.aff}_*(E)$.
